\section{Conceptos fundamentales de la visualización de datos}
\label{fundamentals}
    \subsection{Codificación visual y tipo de variable}
    \label{visual_encoding}
        Toda visualización es una función que asigna valores de datos a atributos
        gráficos que depende del tipo de variable, algunas de ellas son: las variables
        nominales (como el sector del cliente) admiten codificación por posición
        categórica, forma o matiz de color, pero no por longitud continua, las ordinales
        exigen una codificación que preserve el orden y las cuantitativas (consumo en
        kWh, costo en COP) que admiten posición, longitud y área con precisiones muy
        distintas entre sí \cite{castro2026ciencia}.

    \subsection{Jerarquía de precisión de los canales visuales}
    \label{visual_channels}
        El aporte experimental de Cleveland y McGill es la base cuantitativa del diseño
        gráfico moderno, puesto que, al medir el error en la estimación de magnitudes,
        se establece el ordenamiento de precisión decreciente
        \cite{cleveland1984graphical}, el cual se da de la siguiente manera:

        \vspace{10pt}

        \begin{itemize}
            \item Posición sobre una escala común (gráficos de dispersión, barras
            alineadas).

            \item Posición sobre escalas no alineadas (paneles múltiples).

            \item Longitud (barras).

            \item Ángulo y pendiente (gráficos de torta).

            \item Área (burbujas).

            \item Volumen, curvatura y saturación de color.
        \end{itemize}

    \subsection{Correspondencia entre tarea analítica y tipo de gráfico}
    \label{task_chart_correspondence}
        Elegir un tipo de gráfico no es una decisión estética, sino la respuesta a una
        pregunta analítica concreta, ya que no es lo mismo examinar cómo se distribuye
        una variable, que comparar magnitudes entre categorías o buscar la relación
        entre dos medidas, por consiguiente, cada una de estas tareas tiene formas
        gráficas que la responden con mayor precisión que otras, apoyándose en la
        jerarquía de canales visuales descrita en la sección \ref{visual_channels},
        sintetizando dicha correspondencia justificando la elección de la siguiente
        manera:

        \begin{table}[H]
            \centering
            \caption{Correspondencia entre tarea analítica y forma gráfica}
            \begin{center}
                \footnotesize
                \begin{tabular}{|>{\raggedright\arraybackslash}p{1.9cm}|>{\raggedright\arraybackslash}p{1.7cm}|>{\raggedright\arraybackslash}p{2.9cm}|}
                    \hline
                    \textbf{Tarea analítica} & \textbf{Gráfico} & \textbf{Canal y justificación} \\
                    \hline
                    Distribución de una variable & Histograma & Posición y longitud; revela asimetría y multimodalidad \\
                    \hline
                    Comparación entre categorías & Barras & Longitud sobre línea base común desde cero \\
                    \hline
                    Relación entre dos variables & Dispersión & Posición en dos escalas comunes \\
                    \hline
                    Dispersión y atípicos por grupo & Diagrama de caja & Posición; resume cinco estadísticos \\
                    \hline
                    Composición de un total & Barras horizontales ordenadas & Longitud; sustituye al ángulo de la torta \\
                    \hline
                \end{tabular}
                \label{task_chart_mapping}
            \end{center}
        \end{table}